% Chapter content for NTAS 15 - extracted from NTAS15_data_report.tex
% This file contains only the body content without preamble or document structure



\section{Introduction}
This chapter documents the missing-data handling workflow for \projectname{} \casenumber{} and its role in constructing the continuous merged time series.

\subsection{Deployment Summary}
\begin{itemize}
	\item \textbf{Deployment Period:} From \DeploymentStart{} to \DeploymentEnd{}
	\item \textbf{Instruments Used:} \FirstInstrumentType{} \FirstInstrument{} and \SecondInstrumentType{} \SecondInstrument{}
	\item \textbf{Merged Data:} No data were available for the merged dataset.
	\item \textbf{Overlap with Previous Deployment (\projectname{} \previouscasenumber{}):} no
	\item \textbf{Notes:} Due to a configuration error, neither of the SBE-16s deployed on \projectname{}-\casenumber{} returned data. This report describes how the missing data were handled as part of the merged time series.
\end{itemize}

\section{Pre-Processing}
\subsection{Data Loading and Initial Processing}
\begin{itemize}
	\item No raw data files were available for processing.
\end{itemize}

\section{Data Export}
\begin{itemize}
	\item Fill missing data or create non-existent variables with -99999 values and encode them as missing values.
	\item Carry forward time series for all five variables even if not present, by setting missing values to -99999.
	\item Output files:
	\texttt{NTAS15\_\FirstInstrument\_truncated.nc} and \texttt{NTAS15\_\SecondInstrument\_truncated.nc}.
\end{itemize}

\section{Merged Dataset}
\subsection{Metadata and Selection Strategy}
All necessary metadata are incorporated from individual deployments and retained in the merged file.

\subsection{Overlap Analysis}
\begin{itemize}
	\item \textbf{Overlap determination:} No overlap was available between the end of NTAS-14 and the start of NTAS-15.
	\item \textbf{Merging Method:} Data were merged at the truncation time for the NTAS-14 record.
\end{itemize}

